
Synthetic data has become integral to language model training.
Instruction tuning, knowledge distillation, and data augmentation all routinely use model-generated text, and as synthetic corpora are recycled into training pipelines, the process often becomes recursive: each generation of models is fine-tuned on outputs from the previous one.
Recent work demonstrates that such recursive synthetic training triggers \textit{model collapse}, a progressive erosion of distributional tails and downstream capability~\cite{shumailov2024collapse, dohmatob2024tails, alemohammad2024mad}.
The problem is especially pressing as web-scraped corpora increasingly contain model-generated content, a trend that introduces implicit recursive loops even without deliberate data reuse.

While model collapse is now well-characterized through theoretical analysis of tail decay~\cite{dohmatob2024tails} and empirical studies across modalities~\cite{shumailov2024collapse, alemohammad2024mad}, its effect on the geometric structure of learned representations remains poorly understood.
As recursive training erodes distributional diversity, the embedding space should contract onto a lower-dimensional subspace.
Characterizing this spectral trajectory could yield computationally cheap collapse indicators from internal model states alone, without text generation or benchmark evaluation.

Current collapse detection depends on downstream metrics such as perplexity and benchmark accuracy, or on analysis of generated text through lexical diversity measures.
Both approaches require either costly inference or external evaluation infrastructure at every generation.
SIGMA~\cite{gu2026sigma} introduced spectral monitoring through embedding covariance analysis and showed that spectral indicators can track collapse, but validated only on OPT-125M over short generation horizons.
Three gaps remain open.
First, no work has systematically characterized spectral dynamics beyond a few generations at 410M+ parameters.
Second, the temporal relationship between spectral indicators and downstream degradation has not been examined.
Third, no study has connected spectral early warning signals to closed-loop collapse prevention.

We address these gaps by tracking spectral indicators of the embedding Gram matrix across 10 recursive generations of Pythia-410M-deduped~\cite{biderman2023pythia}, with pilot validation at the 1B scale.
Our analysis reveals a reproducible three-phase pattern (Figure~\ref{fig:spectral_trajectory}): spectral decompression (Gen~0--2), where effective rank ($r_{\text{eff}}$) and spectral log-volume (SLV) rise while downstream performance already declines; collapse onset (Gen~2--4), where spectral indicators reverse; and accelerating collapse (Gen~4--10), where contraction proceeds 2.43$\times$ faster than during onset.
Among these indicators, SLV is particularly suited for automated monitoring: it requires only a forward pass on held-out data (no text generation), captures distributed volume changes across all spectral directions, and declines monotonically after its peak.
We further demonstrate a proof-of-concept proportional controller that uses SLV to guide real-synthetic data mixing, and identify a signal attenuation paradox where successful intervention dampens the spectral signal the controller relies on.

Our contributions are:
\begin{itemize}
\item \textbf{C1:} A systematic spectral characterization of model collapse across 10 generations, revealing a three-phase pattern (decompression, onset, acceleration at 2.43$\times$ rate) validated on Pythia-410M and confirmed at 1B scale (\S\ref{subsec:phase0}, \S\ref{subsec:cross_scale}, Table~\ref{tab:phase0}).
\item \textbf{C2:} Identification of SLV as a reliable, computationally efficient collapse monitor with monotonic post-peak decline, together with the HellaSwag-spectral asymmetry showing that downstream degradation precedes spectral collapse (\S\ref{subsec:phase0}, \S\ref{subsec:cross_scale}).
\item \textbf{C3:} A closed-loop SLV-guided mixing controller (proof of concept) and the signal attenuation paradox insight for future controller design (\S\ref{subsec:controller}, \S\ref{subsec:phase1}).
\end{itemize}